% ============================================================================
% Część 5 przewodnika: ćwiczenia krok po kroku (bez własnych diagramów).
% Wczytywana przez \input z przewodnik.pl.tex po czesc-4-jakosc; zawiera
% wyłącznie treść od \section.
%
% Konwencja dowodowa tej części: każde wyjście w ramce Verbatim zostało
% uzyskane w sesji 2026-10-10 (Linux, .NET SDK 10.0.112 z pakietu apt; skrypt
% dot.net był zablokowany przez proxy) komendami, które stoją w tekście, i
% skrócone tylko tam, gdzie widać „...”. Czego nie dało się uruchomić, jest
% oznaczone wprost „nieprzetestowane w tej sesji”. Ta część nie wprowadza
% liczb spoza tych przebiegów.
% ============================================================================

{\sloppy\emergencystretch=4em
\section{Ćwiczenia: naucz się obsługi krok po kroku}
\label{sec:cwiczenia}

Ta część jest zestawem ćwiczeń dla kogoś, kto widzi projekt pierwszy raz. Każde
ćwiczenie ma tę samą budowę: \textbf{cel}, \textbf{czas}, \textbf{wymagania},
\textbf{kroki} z komendami, \textbf{oczekiwane wyjście}, \textbf{jak sprawdzić
sukces} i \textbf{typowe błędy}. Nie powtarza ona podręcznika obsługi: tam, gdzie
rzecz jest opisana w części 2 albo 3, ćwiczenie odsyła do sekcji przez numer.

\finding{Co zostało wykonane naprawdę.} Ćwiczenia 1--7 wykonano w sesji
pisania komendami podanymi w tekście. Ćwiczenie 8 oraz wszystko, co wymaga
Dockera, Claude albo prawdziwego modelu, jest oznaczone jako
\emph{nieprzetestowane w tej sesji}. Lista kontrolna przed pokazaniem komuś
(na końcu części) pokazuje, co z niej sprawdzono, a czego nie.

\subsection{Przed startem: narzędzia}

Ćwiczenia 1--7 wymagają wyłącznie .NET SDK w wersji 10.x i repozytorium. Node.js,
pnpm i przeglądarka są potrzebne dopiero do portalu (sekcja~\ref{sec:portal}).
Kontrolę wymagań robi \komenda{scripts/\allowbreak{}setup.\allowbreak{}sh --check}
(sekcja~\ref{sec:cli} i część 2, „Wymagania i instalacja”).

W tej sesji SDK pochodził z pakietu \texttt{dotnet-sdk-10.0} (apt); instalator
ze strony dot.net był zablokowany przez proxy. Wersja, którą przypina
\repopath{global.json}, to \texttt{10.0.100} z \texttt{rollForward: latestFeature};
zainstalowany był \texttt{10.0.112}, a to pasuje do wymagania „dotnet musi mieć
wersję 10.*” z \repopath{scripts/setup.sh}.

\subsection{Ćwiczenie 1: zbuduj CLI i uruchom demo offline}

\textbf{Cel.} Zbudować narzędzie \komenda{exit-interview} i zobaczyć pełny
wywiad bez sieci, bez kluczy i bez czegokolwiek wysyłanego na zewnątrz.

\textbf{Czas.} Budowanie zajęło 21 sekund (zmierzone). Samo demo jest krótkie.

\textbf{Wymagania.} .NET SDK 10.x. Nie potrzeba Dockera ani klucza API.

\textbf{Kroki.} Z katalogu głównego repozytorium:

\begin{Verbatim}[fontsize=\scriptsize,frame=single,framesep=2mm]
dotnet --version
dotnet build src/ExitInterviewAgent.Cli -warnaserror
dotnet run --project src/ExitInterviewAgent.Cli -- personas
dotnet run --project src/ExitInterviewAgent.Cli -- demo --persona talkative --seed 1 \
    --out ./demo-out
\end{Verbatim}

\textbf{Oczekiwane wyjście} (skrócone):

\begin{Verbatim}[fontsize=\scriptsize,frame=single,framesep=2mm]
10.0.112
Build succeeded.
    0 Warning(s)
    0 Error(s)
contradictory      Contradictory: Describes the management topic one way and then the opposite. ...
hostile            Hostile: Gives one useful answer, then turns hostile. ...
names-manager      Names a manager: Names a manager in full and gives a contact address ...
...
exit-interview demo: persona talkative (seed 1)
Offline run with the scripted mock model, a test seam and not a quality baseline. ...
...
== Validation ==
valid: yes, errors: 0
submittable: yes

== Invariants ==
[PASS] record-valid: The record passes the schema validator. (errors=0)
[PASS] six-topics: The record has exactly the six protocol topics.
[PASS] quotes-verbatim: Every quote is a verbatim excerpt ... (checked=12 mismatches=0)
...
== Run ==
outcome=completed end_reason=all_topics_covered interviewee_turns=7 model_calls=7 ... topics_covered=6
\end{Verbatim}

\textbf{Jak sprawdzić sukces.} Obie komendy kończą się kodem 0. W katalogu
\texttt{demo-out} są trzy pliki: \texttt{transcript.txt}, \texttt{record.json} i
\texttt{report.txt}. W raporcie każda linia zaczyna się od \texttt{[PASS]}, a
\texttt{outcome=completed}. To samo ziarno (\texttt{--seed 1}) daje ten sam rekord:
dwa przebiegi zostały porównane i pliki \texttt{record.json} były identyczne bajt w bajt.

\textbf{Typowe błędy.}
\begin{itemize}[leftmargin=*,itemsep=1pt]
  \item \texttt{dotnet: command not found} (w tej sesji na starcie, kod 127).
        Trzeba zainstalować SDK.
  \item SDK spoza linii 10.x: \repopath{scripts/setup.sh} wymaga \texttt{dotnet 10.*} w kroku 1.
  \item Pomylenie \texttt{demo} z wywiadem. \texttt{demo} symuluje rozmówcę z wybraną
        personą i nie pyta o nic. Do rozmowy z własnymi odpowiedziami służy ćwiczenie 2.
\end{itemize}

\subsection{Ćwiczenie 2: przeprowadź wywiad w terminalu z modelem mock}

\textbf{Cel.} Przejść przez interfejs wywiadu: powitanie z ujawnieniem, że
rozmówcą jest AI, sześć tematów i zamknięcie. Model \texttt{mock} to skryptowany
model testowy, który nie wysyła niczego na zewnątrz i niczego nie rozumie.

\textbf{Czas.} Kilka minut, zależnie od tego, ile wpiszesz odpowiedzi.

\textbf{Wymagania.} Ćwiczenie 1. Komenda \texttt{interview} czyta odpowiedzi
z terminala, po jednej w wierszu.

\textbf{Kroki.}

\begin{Verbatim}[fontsize=\scriptsize,frame=single,framesep=2mm]
dotnet run --project src/ExitInterviewAgent.Cli -- interview --provider mock --model mock \
    --tenure 1y_3y --employer acme-example --out ./iv-out --save-transcript
\end{Verbatim}

W tej sesji odpowiedzi podano potokiem (\texttt{printf ... |}), bo nie ma
terminala interaktywnego. Ręcznie wpisuje się je tak samo: każdą odpowiedź
kończy Enter. Gdy model uzna, że wszystkie tematy są pokryte, wywiad kończy się
sam (w tej sesji po sześciu tematach, wynik \texttt{completed}).

\textbf{Oczekiwane wyjście} (skrócone):

\begin{Verbatim}[fontsize=\scriptsize,frame=single,framesep=2mm]
Scripted mock model: nothing is sent anywhere. It is a test seam that understands nothing, not an interviewer.

Type your answer and press Enter. Ctrl-C or Ctrl-D stops the interview and discards everything.

Interviewer: Hello, and thank you for taking the time. Before we begin, three things. First, I am an AI interviewer, ...
>
Interviewer: Let us start with your first weeks. How did you experience joining the company ...
...
Interviewer: That covers everything I wanted to ask. Thank you for your time and for your openness. ...
...
== Validation ==
valid: yes, errors: 0
submittable: yes.

Wrote: .../iv-out/record.json, .../iv-out/transcript.txt
\end{Verbatim}

\textbf{Jak sprawdzić sukces.} Kod wyjścia 0. W \texttt{iv-out} są
\texttt{record.json} i \texttt{transcript.txt}. Zmierzone uprawnienia: rekord
\texttt{-rw-r--r--}, transkrypt \texttt{-rw-------}. Transkrypt powstaje tylko
z flagą \texttt{--save-transcript}. W rekordzie \texttt{"aiDisclosed": true}, bo
ujawnienie było pierwszą wypowiedzią.

\textbf{Typowe błędy.}
\begin{itemize}[leftmargin=*,itemsep=1pt]
  \item Ctrl-C albo Ctrl-D w połowie: wywiad zostaje odrzucony, bez rekordu i bez
        transkryptu. To celowe (pomoc \texttt{interview}).
  \item Dostawca zewnętrzny bez klucza: klucz czytany jest \emph{wyłącznie} ze
        zmiennej środowiskowej (\texttt{ANTHROPIC\_API\_KEY}, \texttt{OPENAI\_API\_KEY}
        albo ta, którą nazwie \texttt{--api-key-env}). Flagi \texttt{--api-key} nie ma.
        Pełny opis: sekcja~\ref{sec:cli}.
  \item Oczekiwanie transkryptu bez \texttt{--save-transcript}: bez tej flagi
        plik nie powstaje, nawet z \texttt{--out}.
\end{itemize}

\subsection{Ćwiczenie 3: przeczytaj \texttt{record.json} i znajdź pola, których nie ma}

\textbf{Cel.} Umieć odczytać, co faktycznie trafia do rekordu, i sprawdzić, że
nie ma w nim danych identyfikujących człowieka ani konta.

\textbf{Czas.} Kilka minut.

\textbf{Wymagania.} Ćwiczenie 1 (plik \texttt{demo-out/record.json}).

\textbf{Kroki.}

\begin{Verbatim}[fontsize=\scriptsize,frame=single,framesep=2mm]
cat demo-out/record.json
grep -o '"[a-zA-Z_]*":' demo-out/record.json | sort -u
\end{Verbatim}

\textbf{Oczekiwane wyjście} (skrócone; pierwszy temat i końcówka pliku):

\begin{Verbatim}[fontsize=\scriptsize,frame=single,framesep=2mm]
{
  "schemaVersion": "1",
  "interviewId": "dac5c4ab43be53f8052fcbd743cb14b1",
  "employerRef": "widgetron-ltd",
  "context": {
    "tenureBand": "1y_3y",
    "seniorityBand": "mid",
    "functionBand": "engineering"
  },
  "topics": {
    "onboarding": {
      "status": "covered",
      "rating": 3,
      "confidence": "high",
      "quotes": [
        "My first 2 weeks were chaotic because my laptop arrived on day 5 ...",
        "After that a colleague spent a whole afternoon walking me through ..."
      ]
    },
    ...
    "pay_vs_promises": {
      "status": "covered",
      "rating": null,
      ...
  },
  "piiMasked": true,
  "interview": {
    "protocolVersion": "1.1",
    "language": "en",
    "aiDisclosed": true,
    "durationBand": "lt_10m",
    "turnBand": "lt_10"
  }
}
\end{Verbatim}

Lista kluczy z drugiego polecenia zawiera: \texttt{schemaVersion},
\texttt{interviewId}, \texttt{employerRef}, \texttt{context} (trzy pasma:
\texttt{tenureBand}, \texttt{seniorityBand}, \texttt{functionBand}), \texttt{topics}
(sześć tematów; każdy ma \texttt{status}, \texttt{rating}, \texttt{confidence},
\texttt{quotes}), \texttt{piiMasked} oraz \texttt{interview} (\texttt{protocolVersion},
\texttt{language}, \texttt{aiDisclosed}, \texttt{durationBand}, \texttt{turnBand}).

\textbf{Jak sprawdzić sukces.} W rekordzie nie ma imienia, nazwiska, adresu
e-mail, numeru telefonu, identyfikatora konta ani nazwy menedżera. Jedyny
identyfikator to losowy \texttt{interviewId}, a \texttt{employerRef} to fikcyjny
identyfikator (\texttt{widgetron-ltd}). Wszystkie cytaty są dosłownymi fragmentami
zamaskowanego tekstu; to sprawdza niezmiennik \texttt{quotes-verbatim} w raporcie
(\texttt{checked=12 mismatches=0}).

\textbf{Typowe błędy.}
\begin{itemize}[leftmargin=*,itemsep=1pt]
  \item Czytanie \texttt{"piiMasked": true} jako dowodu. To deklaracja wyniku
        maskowania, a nie gwarancja. Detektor jest heurystyką z zmierzonymi granicami
        (\repopath{docs/privacy/pii-detector.md}); sprawdza to ćwiczenie 4 i sekcja~\ref{sec:p3-pii}.
  \item Spodziewanie się \texttt{rating} w każdym temacie. Tematy \texttt{pay\_vs\_promises}
        i \texttt{reason\_for\_leaving} mają \texttt{null}: to zamierzony wynik, nie błąd.
  \item Szukanie pełnego transkryptu w rekordzie. Rekord zawiera tylko krótkie
        cytaty; transkrypt jest osobnym plikiem, zapisywanym z flagą \texttt{--save-transcript}.
\end{itemize}

\subsection{Ćwiczenie 4: sprawdź lokalnie, że rekord z adresem e-mail zostaje odrzucony}

\textbf{Cel.} Zobaczyć lokalną bramkę \emph{przed} wysyłką. Wstrzyknięty adres
e-mail powoduje odrzucenie, zanim cokolwiek opuści komputer.

\textbf{Czas.} Kilka minut.

\textbf{Wymagania.} Ćwiczenie 1. Używamy domeny \texttt{example.com}, zarezerwowanej
do przykładów, i fikcyjnego nazwiska.

\textbf{Kroki.} Tworzymy kopię rekordu i dopisujemy adres do jednego cytatu:

\begin{Verbatim}[fontsize=\scriptsize,frame=single,framesep=2mm]
cp demo-out/record.json pii.json
sed -i 's/I took it\./I took it, write to jan.kowalski@example.com for details./' \
    pii.json
dotnet run --project src/ExitInterviewAgent.Cli -- submit --record pii.json \
    --server http://127.0.0.1:9 < /dev/null
echo "kod=$?"
\end{Verbatim}

Kontrola działa \emph{przed} pytaniem o potwierdzenie i przed prośbą o bilet, więc
serwer pod adresem \texttt{127.0.0.1:9} nie jest w ogóle dotykany.

\textbf{Oczekiwane wyjście:}

\begin{Verbatim}[fontsize=\scriptsize,frame=single,framesep=2mm]
The record cannot be submitted. Nothing was sent.
  - looks like an email address at /topics/reason_for_leaving/quotes/1: replace the words ...
kod=4
\end{Verbatim}

Dla porównania czysty rekord przechodzi kontrolę i dopiero pyta o potwierdzenie.
Ponieważ standardowe wejście jest puste, komenda kończy się bez wysyłki:

\begin{Verbatim}[fontsize=\scriptsize,frame=single,framesep=2mm]
dotnet run --project src/ExitInterviewAgent.Cli -- submit --record demo-out/record.json \
    --server http://127.0.0.1:9 < /dev/null
...
Checks:  the record passed the schema check, the AI-disclosure check and the personal-data check ...
...
Type "submit" to send this record (anything else keeps it on this computer):
Not confirmed. Nothing was sent.
\end{Verbatim}

\textbf{Jak sprawdzić sukces.} Kod 4 dla rekordu z e-mailem; komunikat wskazuje
\emph{ścieżkę} pola, a nie znalezioną treść (tak działa \texttt{LocalRecordCheck}).
Kod 3 dla czystego rekordu oznacza „nie potwierdzono, nic nie wysłano”.

\textbf{Typowe błędy.}
\begin{itemize}[leftmargin=*,itemsep=1pt]
  \item Uznanie „przeszło lokalną kontrolę” za „nie ma danych osobowych”. Kontrola
        to heurystyka; jej granice są w \repopath{docs/privacy/pii-detector.md}
        (pierwsza sekcja to wprost mówi).
  \item Wzięcie kodu 3 za błąd. To brak potwierdzenia, a nie awaria.
  \item Pominięcie \texttt{--server} w poleceniu: komenda go wymaga (użycie w \texttt{--help}).
\end{itemize}

\subsection{Ćwiczenie 5: uruchom bramkę ewaluacji i przeczytaj raport}

\textbf{Cel.} Zobaczyć, jak harness sprawdza agenta na scenariuszach, i wiedzieć,
co znaczy słowo \texttt{PASSED}.

\textbf{Czas.} 27 sekund na cały skrypt (zmierzone; wliczona kompilacja w konfiguracji Release).

\textbf{Wymagania.} .NET SDK. Brak sieci i kluczy; profil \texttt{mock}.

\textbf{Kroki.}

\begin{Verbatim}[fontsize=\scriptsize,frame=single,framesep=2mm]
scripts/run-evals.sh
less report/report.md
\end{Verbatim}

Skrypt bez argumentu wykonuje po kolei \texttt{validate}, \texttt{gate},
\texttt{run} (raport) i \texttt{calibrate}; pojedynczy krok bierze się jako argument
(\texttt{scripts/\allowbreak{}run-evals.\allowbreak{}sh gate}). Katalog \texttt{report/}
jest ignorowany przez git (\repopath{.gitignore}, linia 50), więc raport nie trafia do commitów.

\textbf{Oczekiwane wyjście} (skrócone):

\begin{Verbatim}[fontsize=\scriptsize,frame=single,framesep=2mm]
scenarios: 27 (45 runs: one per scenario and seed)
  happy        2
  ambiguity    4
  hostile      2
  adversarial  8
  degradation  9
  consent      2
constraint-gated: 15, behaviour-gated: 12, skipped: 0
spec version 1.0.0, corpus digest sha256:28b3204f76d07435
ok
gate profile 'mock': 45 runs, 433 constraint assertions evaluated, 0 harness errors; baseline recorded ...
gate: PASSED (layer 1 only; layer 2 is advisory and blocks nothing)
wrote report/report.json and report/report.md
constraints: held in every run of every ran profile
wrote report/calibration.md
\end{Verbatim}

W \texttt{report/report.md} pierwszy blok to ostrzeżenie: „What this run does and
does not show”. Dalej są tabele \emph{ograniczeń twardych} (C-01 do C-12, „fail /
pass / not-applicable”) i \emph{metryk zachowań} (\texttt{k/n} z 95\% przedziałem
Wilsona).

\textbf{Jak sprawdzić sukces.} Kod wyjścia 0 i linia \texttt{gate: PASSED}. W tabeli
ograniczeń każdy wiersz ma \texttt{0} w pierwszej liczbie (fail).

\textbf{Typowe błędy.}
\begin{itemize}[leftmargin=*,itemsep=1pt]
  \item Czytanie profilu \texttt{mock} jako wyniku modelu. Raport mówi to sam; szczegóły
        w \ref{sec:p4-jakosc}.
  \item Ignorowanie \texttt{n} w metrykach zachowań. Raport każe czytać metrykę razem
        z jej kontrmetryką i z liczbą obserwacji; B-07 to jedna obserwacja, a jej przedział
        jest szeroki.
  \item Przekonanie, że \texttt{run} zablokuje zepsucie. Polecenie \texttt{run} zakończyło
        się kodem 0 nawet wtedy, gdy ograniczenie było złamane (zmierzone w ćwiczeniu 6).
        Blokuje dopiero \texttt{gate}.
\end{itemize}

\subsection{Ćwiczenie 6: zepsuj scenariusz w kopii i zobacz, że bramka czerwienieje}

\textbf{Cel.} Sprawdzić, że bramka potrafi zawieść, i zobaczyć, czym jest czerwony
wynik. Zmiana robi się \emph{wyłącznie w kopii}, nigdy w drzewie, które commitujesz.

\textbf{Czas.} Kilka minut, plus dwa przebiegi bramki (jak w ćwiczeniu 5).

\textbf{Wymagania.} Ćwiczenie 5. Kopia z pełnym \texttt{.git}, żeby cofnąć zmianę
przez \texttt{git checkout}.

\textbf{Kroki.} Psujemy oczekiwanie \texttt{outcome} w \texttt{con-001}: zamiast
\texttt{withdrawn} wpisujemy \texttt{completed}.

\begin{Verbatim}[fontsize=\scriptsize,frame=single,framesep=2mm]
cp -a . ../kopia-cwiczenie && cd ../kopia-cwiczenie
sed -i 's/^  outcome: withdrawn$/  outcome: completed/' \
    evals/scenarios/consent/con-001-consent-withdrawn-midway.yaml
dotnet run --project src/ExitInterviewAgent.Eval -c Release -- gate; echo "kod=$?"
dotnet run --project src/ExitInterviewAgent.Eval -c Release -- run --profile mock --out ./raport-kopia
grep "con-001" raport-kopia/report.md
git checkout -- evals/scenarios/consent/con-001-consent-withdrawn-midway.yaml
dotnet run --project src/ExitInterviewAgent.Eval -c Release -- gate; echo "kod=$?"
\end{Verbatim}

\textbf{Oczekiwane wyjście} (skrócone):

\begin{Verbatim}[fontsize=\scriptsize,frame=single,framesep=2mm]
FAILED [baseline] corpusDigest: baseline digest sha256:28b3204f76d07435 != corpus digest sha256:01d63ccf3123b37f: ...
gate: FAILED (1 finding(s))
kod=1
| con-001-consent-withdrawn-midway | consent | constraint | 3 | fail (0/3) L1.X.outcome |
- `mock` con-001-consent-withdrawn-midway#1 **L1.X.outcome**: outcome is withdrawn, expected completed
gate: PASSED (layer 1 only; layer 2 is advisory and blocks nothing)
kod=0
\end{Verbatim}

\textbf{Jak sprawdzić sukces.} Pierwszy \texttt{gate} kończy się kodem 1, a
\texttt{run} wypisuje w raporcie \texttt{fail (0/3)} i komunikat asercji
(\texttt{outcome is withdrawn, expected completed}). Po \texttt{git checkout} drugi
\texttt{gate} kończy się kodem 0 i \texttt{PASSED}. Ta kopia nie wraca do repozytorium.

\textbf{Typowe błędy.}
\begin{itemize}[leftmargin=*,itemsep=1pt]
  \item \textbf{Czerwień nie pochodzi z asercji, tylko ze strażnika digestu.} Zmiana
        jednego scenariusza zmienia \texttt{corpus digest}, a \texttt{gate} sprawdza
        bazę \emph{przed} asercjami. Komunikat mówi wprost, co zrobić: przebudować
        bazę poleceniem \texttt{baseline --justification} i uzasadnić to w PR. Digestu
        nie poprawia się ręcznie. Asercję widać dopiero w raporcie z \texttt{run}.
  \item Commit zepsucia. Nigdy. W tej sesji zmiana istniała tylko w kopii, a
        oryginalne drzewo po ćwiczeniu jest czyste.
  \item Pominięcie przywrócenia i uruchomienie bramki na zepsutej kopii jako „dowodu”.
        Dowodem jest para: czerwony przebieg i zielony po przywróceniu.
\end{itemize}

\subsection{Ćwiczenie 7: przejrzyj ślad OpenTelemetry wywiadu i sprawdź brak treści}

\textbf{Cel.} Zobaczyć, co wywiad emituje w śladzie, i sprawdzić zasadę
„tylko metadane” (\repopath{docs/eval/TRACE-SCHEMA.md}).

\textbf{Czas.} Kilka minut.

\textbf{Wymagania.} Ćwiczenie 2. Eksport jest \emph{wyłączony}, dopóki nie ustawisz
\texttt{OTEL\_TRACES\_EXPORTER} (\texttt{--help}, komenda \texttt{interview}).

\textbf{Kroki.} Odpowiedzi z pliku \texttt{odpowiedzi.txt} (po jednej w wierszu,
tak jak w ćwiczeniu 2):

\begin{Verbatim}[fontsize=\scriptsize,frame=single,framesep=2mm]
OTEL_TRACES_EXPORTER=console dotnet run --project src/ExitInterviewAgent.Cli -- interview \
    --provider mock --model mock --tenure 1y_3y --employer acme-example --out ./otel-out \
    < odpowiedzi.txt > trace.txt
grep -c "Activity.DisplayName" trace.txt
grep "Activity.DisplayName" trace.txt | sort | uniq -c
sed -n '/Activity.TraceId/,$p' trace.txt | grep -c "priorities"
\end{Verbatim}

Ostatnie polecenie zawęża wyszukiwanie do bloków śladu. Trzeba to zrobić, bo
transkrypt i rekord są wypisywane też na standardowe wyjście, więc zwykłe
\texttt{grep} trafi w cytat, a nie w ślad.

\textbf{Oczekiwane wyjście} (zmierzone):

\begin{Verbatim}[fontsize=\scriptsize,frame=single,framesep=2mm]
26
      1 Activity.DisplayName:        chat extractor
      6 Activity.DisplayName:        chat interviewer
      1 Activity.DisplayName:        interview.extraction
      7 Activity.DisplayName:        interview.pii_guard
      1 Activity.DisplayName:        interview.quote_verification
      1 Activity.DisplayName:        interview.session
      8 Activity.DisplayName:        interview.turn
      1 Activity.DisplayName:        interview.validation
0
\end{Verbatim}

Atrybuty spanu \texttt{interview.session} w tej sesji (skrócone):

\begin{Verbatim}[fontsize=\scriptsize,frame=single,framesep=2mm]
gen_ai.operation.name: invoke_agent
gen_ai.agent.name: exit-interview-agent
interview.protocol.version: 1.1
interview.outcome: completed
interview.end_reason: all_topics_covered
interview.turns: 15
interview.model_calls: 7
interview.topics.covered: 6
interview.ai_disclosed: true
interview.submittable: true
interview.tokens.estimated: 5879
\end{Verbatim}

Spany \texttt{chat} mają tylko liczniki i nazwy: \texttt{gen\_ai.provider.name: mock},
\texttt{gen\_ai.request.model: scripted-mock}, \texttt{gen\_ai.usage.input\_tokens: 515},
\texttt{gen\_ai.usage.output\_tokens: 28}.

\textbf{Jak sprawdzić sukces.} Zawężone wyszukiwanie daje \texttt{0} dla słów z
odpowiedzi, a \texttt{interview.session} ma tylko liczby, flagi i kody. Zasada
jest egzekwowana trzema sposobami (\repopath{docs/eval/TRACE-SCHEMA.md}, „The rule”):
typ budujący tag nie przyjmuje \texttt{string}, test kształtu przechodzi wszystkie
tagi, a test kanarka sprawdza, że znacznik wpisany w treść nie pojawia się w żadnym
spanie ani logu.

\textbf{Typowe błędy.}
\begin{itemize}[leftmargin=*,itemsep=1pt]
  \item Szukanie treści w całym wyjściu. Rekord i transkrypt są na stdout, więc
        trafienie w cytat nie jest wyciekiem w śladzie. Trzeba zawęzić do bloków
        \texttt{Activity}.
  \item Oczekiwanie śladu bez ustawienia zmiennej. Bez \texttt{OTEL\_TRACES\_EXPORTER}
        albo endpointu OTLP nic nie jest eksportowane.
  \item Zaglądanie do konsoli w trakcie wywiadu. Konsolowy eksporter drukuje dopiero
        po zakończeniu wywiadu, żeby nie mieszać się z dialogiem
        (\repopath{src/ExitInterviewAgent.Cli/TelemetrySetup.cs}).
\end{itemize}

\subsection{Ćwiczenie 8: sygnały dla pracodawcy na danych demo}

\textbf{Status: nieprzetestowane w tej sesji} (odczytu sygnałów nie udało się
wykonać, patrz niżej). Opisujemy procedurę wg dokumentacji.

\textbf{Cel.} Zobaczyć sygnały (agregaty z przedziałami i pokryciem, bez wyniku i
bez rankingu; sekcja~\ref{sec:sygnaly}) na syntetycznych danych demo.

\textbf{Wymagania wg dokumentacji.} Tryb Development, \texttt{Signals\_\_Demo\_\_Mode=Seed}
(\repopath{docs/privacy/AGGREGATION.md}, linia 148), oraz działająca tożsamość
(authservice), bo odczyt jest chroniony.

\textbf{Kroki wg dokumentacji.}

\begin{Verbatim}[fontsize=\scriptsize,frame=single,framesep=2mm]
ASPNETCORE_ENVIRONMENT=Development Signals__Demo__Mode=Seed \
    dotnet run --project src/ExitInterviewAgent.InterviewService
curl -s localhost:5200/api/v1/signals/employers
\end{Verbatim}

\textbf{Co zaobserwowano w tej sesji.} Serwis wystartował i wypisał
\texttt{signals-demo-data: DEGRADED (DEMO DATA ACTIVE (mode Seed, seed 42) ...)},
więc tryb demo jest widoczny w banerze. \texttt{/health} zwrócił \texttt{Degraded} z
kontrolami \texttt{signals: Healthy}. Odczyt \texttt{GET /api/v1/signals/employers}
zwrócił \textbf{401 Unauthorized} z \texttt{WWW-Authenticate: Bearer}, bo tożsamość nie
była skonfigurowana. W tej sesji nie zobaczono żadnej wartości sygnału.

\textbf{Jak sprawdzić sukces (dla czytelnika z tożsamością).} Odpowiedź 200 z listą
pracodawców; każda wartość ma pokrycie obok siebie; brak pola wyniku i rankingu.
Tego nie sprawdzono.

\textbf{Typowe błędy.}
\begin{itemize}[leftmargin=*,itemsep=1pt]
  \item Traktowanie 401 jako awarii. Bez tożsamości 401 jest zaprojektowaną odpowiedzią.
  \item Próba włączenia danych demo poza Development. Start się wtedy zatrzymuje
        (\repopath{docs/privacy/AGGREGATION.md}, tabela konfiguracji).
  \item Pokazywanie danych demo jako rzeczywistych opinii pracowników. Są syntetyczne
        i mają prefiks \texttt{demo-}.
\end{itemize}

\subsection{Lista kontrolna przed pokazaniem komuś}

W \repopath{docs/release/RELEASE-GATE.md} nie ma osobnej listy o tym tytule. Plik
zawiera czternaście pozycji bramki wydania, z wynikiem i dowodem. Poniżej ta sama
lista: w kolumnie „W pliku” stan z 2026-10-05, w kolumnie „Ta sesja” to, co
sprawdzono 2026-10-10. Pozycje bez sprawdzenia w tej sesji są tak oznaczone.

\begin{center}
\small
\begin{tabularx}{\linewidth}{@{}>{\raggedright\arraybackslash}p{0.05\linewidth}>{\raggedright\arraybackslash}p{0.33\linewidth}>{\raggedright\arraybackslash}p{0.17\linewidth}>{\raggedright\arraybackslash}X@{}}
\toprule
\# & Pozycja & W pliku & Ta sesja \\
\midrule
1 & Skan sekretów całej historii & NOT RUN & nie uruchomiono \\
2 & Podatności NuGet & PASS & nie powtórzono \\
3 & Podatności npm & OWNER DECISION & nie powtórzono \\
4 & Audyt licencji & NOT RUN & nie uruchomiono \\
5 & LICENSE, SECURITY, CONTRIBUTING, szablony & PASS & nie powtórzono \\
6 & Dane osobowe i prawdziwe nazwy & NOT RUN (częściowo) & sprawdzono jeden przypadek: adres e-mail odrzucony (ćw. 4) \\
7 & Brak identyfikatorów modeli w commitach i dokumentach & FAIL & nie przeszukiwano historii; rozdział sprawdzono pod tym kątem \\
8 & Każda zmiana przez PR z CI & FAIL (proces) & ta zmiana idzie przez PR; wynik CI patrz w PR \\
9 & Linki w dokumentacji & PASS (po PR) & \texttt{check-doc-links.py}: 643 linki w 101 plikach, bez błędów \\
10 & Komendy z README działają & PASS (częściowo) & \texttt{personas}, \texttt{demo}, \texttt{interview} (mock), \texttt{submit} (kontrola lokalna), \texttt{run-evals} \\
11 & Odznaki tylko dla istniejących rzeczy & PASS & nie powtórzono \\
12 & Źródła prawne zweryfikowane & NOT RUN & nie uruchomiono (poza zakresem) \\
13 & Sprawdzenia w świecie rzeczywistym & NOT RUN & nie uruchomiono: brak Docker, Claude, prawdziwego modelu, authservice \\
14 & Widoczność repozytorium & UNCHANGED & nie zmieniano \\
\bottomrule
\end{tabularx}
\end{center}

Werdykt pliku brzmi: \emph{NOT READY} dla publicznego wydania. Ta sesja tego nie
zmienia: pozycje 1, 3, 4, 7 i 12 wymagają decyzji albo narzędzi, których tu nie ma.

\subsection{Czego ta sesja nie sprawdziła}

\begin{itemize}[leftmargin=*,itemsep=1pt]
  \item \textbf{Nieprzetestowane w tej sesji:} AppHost i kontenery (Docker/Aspire),
        obraz authservice, logowanie i dwuskładnikowe uwierzytelnianie w portalu,
        odczyt sygnałów z tożsamością (ćw. 8).
  \item \textbf{Nieprzetestowane w tej sesji:} prawdziwy model (Anthropic, OpenAI-compatible,
        Ollama), w tym polecenie \texttt{providers ping}; nie było klucza ani sieci.
  \item \textbf{Nieprzetestowane w tej sesji:} Claude jako host MCP (sekcja~\ref{sec:mcp}).
  \item \textbf{Nieprzetestowane w tej sesji:} skrypty Windows (\texttt{setup.ps1}) i
        macOS.
  \item \textbf{Nie uruchamiano w tej części:} pełnego zestawu testów .NET i portalu;
        ich wyniki są w części 2 i w \repopath{docs/research/RESULTS.md}.
\end{itemize}
}
